\documentclass[aspectratio=169,12pt]{beamer}

\usepackage[T1]{fontenc}
\usepackage[utf8]{inputenc}
\usepackage[brazil]{babel}
\usepackage{amsmath,amssymb,amsfonts}
\usepackage{mathtools}
\usepackage{booktabs}
\usepackage{array}
\usepackage{xcolor}
\usepackage{tikz}
\usepackage{pgfplots}
\pgfplotsset{compat=1.18}
\usetikzlibrary{arrows.meta,positioning,shapes.geometric,decorations.pathreplacing,calc,patterns}

%==============================================================
% TEMA E CORES UFAM
%==============================================================
\usetheme{Madrid}
\usecolortheme{default}

\definecolor{ufamblue}{RGB}{0,51,102}
\definecolor{ufamlightblue}{RGB}{0,102,170}
\definecolor{ufamgold}{RGB}{204,153,0}
\definecolor{ufamgreen}{RGB}{0,102,51}
\definecolor{matrixbg}{RGB}{240,248,255}
\definecolor{resultcolor}{RGB}{0,80,40}
\definecolor{alertred}{RGB}{180,30,30}

\setbeamercolor{palette primary}{bg=ufamblue, fg=white}
\setbeamercolor{palette secondary}{bg=ufamlightblue, fg=white}
\setbeamercolor{palette tertiary}{bg=ufamgold, fg=black}
\setbeamercolor{palette quaternary}{bg=ufamblue, fg=white}
\setbeamercolor{structure}{fg=ufamblue}
\setbeamercolor{section in toc}{fg=ufamblue}
\setbeamercolor{frametitle}{bg=ufamblue, fg=white}
\setbeamercolor{title}{bg=ufamblue, fg=white}
\setbeamercolor{block title}{bg=ufamlightblue, fg=white}
\setbeamercolor{block body}{bg=matrixbg, fg=black}
\setbeamercolor{block title alerted}{bg=alertred, fg=white}
\setbeamercolor{block body alerted}{bg=red!5, fg=black}
\setbeamercolor{block title example}{bg=ufamgreen, fg=white}
\setbeamercolor{block body example}{bg=green!5, fg=black}

\setbeamerfont{frametitle}{size=\large, series=\bfseries}
\setbeamerfont{title}{size=\Large, series=\bfseries}

\setbeamertemplate{navigation symbols}{}
\setbeamertemplate{footline}{
  \leavevmode
  \hbox{
    \begin{beamercolorbox}[wd=.33\paperwidth,ht=2.5ex,dp=1.5ex,center]{palette primary}
      \usebeamerfont{author in head/foot}\insertshortauthor
    \end{beamercolorbox}
    \begin{beamercolorbox}[wd=.34\paperwidth,ht=2.5ex,dp=1.5ex,center]{palette secondary}
      \usebeamerfont{title in head/foot}\insertshorttitle
    \end{beamercolorbox}
    \begin{beamercolorbox}[wd=.33\paperwidth,ht=2.5ex,dp=1.5ex,center]{palette primary}
      \usebeamerfont{date in head/foot}
      \insertframenumber\,/\,\inserttotalframenumber
    \end{beamercolorbox}
  }
}

%==============================================================
\title[Segmenta\c{c}\~ao Morfol\'ogica]{Improved Image Segmentation Method\\Based on Morphological Reconstruction}
\subtitle{Wu, Y., Peng, X., Ruan, K. \textit{et al.} --- \textit{Multimedia Tools and Applications}, 2017}
\author[M.\ Uch\^oa]{Matheus Serr\~ao Uch\^oa}
\institute[UFAM/IComp]{%
  Universidade Federal do Amazonas --- IComp\\
  Processamento Digital de Imagens}
\date{\today}

\begin{document}

%==============================================================
% SLIDE 1 — CAPA
%==============================================================
\begin{frame}[plain]
\titlepage
\end{frame}

%==============================================================
% SLIDE 2 — ROTEIRO
%==============================================================
\begin{frame}{Roteiro}
\tableofcontents
\end{frame}

%==============================================================
\section{Introdu\c{c}\~ao e Motiva\c{c}\~ao}
%==============================================================

%==============================================================
% SLIDE 3
%==============================================================
\begin{frame}{O que \'e a Transformada Watershed?}

\begin{columns}[T]
\begin{column}{0.55\textwidth}
\begin{itemize}
  \item Interpreta a imagem como \textbf{superf\'icie topogr\'afica}
  \item Pixels claros $\to$ ``montanhas''; escuros $\to$ ``vales''
  \item \'Agua sobe simultaneamente a partir de cada m\'inimo local
  \item Onde as \'aguas se encontram $\to$ \textbf{linha de watershed} (fronteira)
  \item Cada regi\~ao inundada $=$ \textbf{bacia de capta\c{c}\~ao} (segmento)
\end{itemize}
\end{column}
\begin{column}{0.42\textwidth}
\centering
\begin{tikzpicture}[scale=0.65]
  % Topographic surface
  \draw[thick, ufamblue, fill=blue!8]
    (0,0) -- (0.5,1.2) -- (1.2,0.3) -- (1.8,1.8) -- (2.5,0.5) --
    (3.2,2.0) -- (3.8,0.4) -- (4.5,1.5) -- (5,0) -- cycle;
  % Water levels
  \fill[blue!30, opacity=0.5] (0,0) -- (0.5,1.0) -- (1.2,0.3) -- (1.2,0) -- cycle;
  \fill[red!30, opacity=0.5] (1.2,0) -- (1.2,0.3) -- (2.5,0.5) -- (2.5,0) -- cycle;
  \fill[green!30, opacity=0.5] (2.5,0) -- (2.5,0.5) -- (3.8,0.4) -- (3.8,0) -- cycle;
  \fill[orange!30, opacity=0.5] (3.8,0) -- (3.8,0.4) -- (5,0) -- cycle;
  % Labels
  \draw[<-, thick, alertred] (1.2,0.3) -- (1.2,-0.5) node[below, font=\tiny]{Watershed};
  \draw[<-, thick, alertred] (2.5,0.5) -- (2.5,-0.5) node[below, font=\tiny]{Watershed};
  \draw[<-, thick, alertred] (3.8,0.4) -- (3.8,-0.5) node[below, font=\tiny]{Watershed};
  \node[font=\tiny, blue!70!black] at (0.6,-0.2) {$B_1$};
  \node[font=\tiny, red!70!black] at (1.85,-0.2) {$B_2$};
  \node[font=\tiny, green!60!black] at (3.15,-0.2) {$B_3$};
  \node[font=\tiny, orange!70!black] at (4.4,-0.2) {$B_4$};
\end{tikzpicture}

\vspace{0.3cm}
{\footnotesize Met\'afora topogr\'afica da watershed.}
\end{column}
\end{columns}

\end{frame}

%==============================================================
% SLIDE 4
%==============================================================
\begin{frame}{O Problema: Sobre-Segmenta\c{c}\~ao}

\begin{alertblock}{Problema Central}
Aplicada diretamente ao gradiente, a watershed trata \textbf{cada m\'inimo local} como uma bacia separada $\Rightarrow$ centenas de regi\~oes esp\'urias.
\end{alertblock}

\vspace{0.3cm}

\begin{columns}[T]
\begin{column}{0.48\textwidth}
\textbf{Causas da sobre-segmenta\c{c}\~ao:}
\begin{itemize}
  \item Ru\'ido na imagem
  \item Varia\c{c}\~oes de textura
  \item Artefatos de quantiza\c{c}\~ao
  \item M\'inimos locais esp\'urios no gradiente
\end{itemize}
\end{column}
\begin{column}{0.48\textwidth}
\textbf{Tentativas anteriores:}
\begin{itemize}
  \item Limiar global \'unico para suprimir m\'inimos rasos
  \item[\color{alertred}$\times$] Falha em imagens com objetos de \textbf{tamanhos diferentes}
  \item[\color{alertred}$\times$] Limiar alto $\to$ perde objetos pequenos
  \item[\color{alertred}$\times$] Limiar baixo $\to$ fragmenta objetos grandes
\end{itemize}
\end{column}
\end{columns}

\end{frame}

%==============================================================
% SLIDE 5
%==============================================================
\begin{frame}{Limita\c{c}\~ao do Limiar \'Unico}

\begin{center}
\begin{tikzpicture}[scale=0.85, >=Stealth]
  % Left scenario
  \node[font=\small\bfseries, ufamblue] at (2.3,3.8) {Limiar Alto};
  \draw[thick] (0,0) rectangle (4.5,3.2);
  % Large object - well segmented
  \draw[thick, ufamgreen, fill=green!10] (0.5,0.5) rectangle (3.5,2.5);
  \node[font=\tiny] at (2.0,1.5) {Objeto grande OK};
  % Small objects - lost
  \draw[thick, red, dashed] (4.0,2.5) circle (0.3);
  \node[font=\tiny, alertred] at (4.0,1.8) {Perdido!};

  % Right scenario
  \begin{scope}[xshift=6.5cm]
  \node[font=\small\bfseries, ufamblue] at (2.3,3.8) {Limiar Baixo};
  \draw[thick] (0,0) rectangle (4.5,3.2);
  % Large object - over-segmented
  \draw[thick, red, dashed] (0.5,0.5) rectangle (3.5,2.5);
  \draw[red, dashed] (2.0,0.5) -- (2.0,2.5);
  \draw[red, dashed] (0.5,1.5) -- (3.5,1.5);
  \node[font=\tiny, alertred] at (2.0,1.0) {Fragmentado!};
  % Small objects - preserved
  \draw[thick, ufamgreen, fill=green!10] (4.0,2.5) circle (0.3);
  \node[font=\tiny] at (4.0,1.8) {OK};
  \end{scope}

  % Arrow between
  \draw[<->, very thick, ufamgold] (5.0,1.6) -- (6.0,1.6);
  \node[font=\tiny\bfseries, ufamgold] at (5.5,2.0) {Trade-off};
\end{tikzpicture}
\end{center}

\vspace{0.3cm}

\begin{block}{Tese Central do Artigo}
\textbf{Nenhum limiar \'unico} resolve simultaneamente objetos de m\'ultiplas escalas.\\
$\Rightarrow$ \'E necess\'aria uma abordagem \textbf{multi-escala} com \textbf{reconstru\c{c}\~ao morfol\'ogica}.
\end{block}

\end{frame}

%==============================================================
\section{Fundamentos de Morfologia Matem\'atica}
%==============================================================

%==============================================================
% SLIDE 6
%==============================================================
\begin{frame}{Opera\c{c}\~oes Morfol\'ogicas B\'asicas}

Sejam $F$ uma imagem e $B$ um elemento estruturante:

\begin{columns}[T]
\begin{column}{0.48\textwidth}
\begin{block}{Dilata\c{c}\~ao $\delta_B$}
\[
  (\delta_B F)(x,y) = \max_{(u,v)\in B} F(x-u,\,y-v)
\]
Expande regi\~oes claras, preenche buracos escuros.
\end{block}

\begin{block}{Abertura $\gamma_B$}
\[
  \gamma_B(F) = \delta_B(\varepsilon_B(F))
\]
Remove picos claros pequenos.
\end{block}
\end{column}
\begin{column}{0.48\textwidth}
\begin{block}{Eros\~ao $\varepsilon_B$}
\[
  (\varepsilon_B F)(x,y) = \min_{(u,v)\in B} F(x+u,\,y+v)
\]
Encolhe regi\~oes claras, elimina pontos claros.
\end{block}

\begin{block}{Fechamento $\varphi_B$}
\[
  \varphi_B(F) = \varepsilon_B(\delta_B(F))
\]
Remove vales escuros pequenos.
\end{block}
\end{column}
\end{columns}

\vspace{0.4cm}

\begin{exampleblock}{Gradiente Morfol\'ogico}
\[
  g(I) = \delta_B(I) - \varepsilon_B(I)
\]
Valores altos nas bordas, $\approx 0$ em regi\~oes homog\^eneas --- \'e a superf\'icie para a watershed.
\end{exampleblock}

\end{frame}

%==============================================================
% SLIDE 7
%==============================================================
\begin{frame}{Reconstru\c{c}\~ao Morfol\'ogica}

\begin{block}{Defini\c{c}\~ao}
Dados um \textbf{marcador} $J$ e uma \textbf{m\'ascara} $I$ (com $J \leq I$ pixel a pixel):
\[
  R^{\delta}_I(J):\quad
  J^{(1)} = \delta_B(J) \wedge I,\quad
  J^{(2)} = \delta_B(J^{(1)}) \wedge I,\quad
  \ldots\quad \text{at\'e } J^{(n+1)} = J^{(n)}
\]
\end{block}

\begin{columns}[T]
\begin{column}{0.45\textwidth}
\centering
\begin{tikzpicture}[scale=0.6]
  % Mask
  \draw[thick, ufamblue, fill=blue!5] plot[smooth] coordinates {(0,0) (0.5,2) (1.5,3) (2.5,1) (3,2.5) (4,3.5) (5,1) (5.5,0)};
  \node[font=\tiny, ufamblue] at (4.5,3.8) {M\'ascara $I$};
  % Marker (eroded)
  \draw[thick, alertred, dashed] plot[smooth] coordinates {(0,0) (0.5,0.5) (1.5,1.5) (2.5,0.3) (3,1.0) (4,2.0) (5,0.3) (5.5,0)};
  \node[font=\tiny, alertred] at (1.0,0.2) {Marcador $J$};
  % Reconstruction result
  \draw[ultra thick, ufamgreen] plot[smooth] coordinates {(0,0) (0.5,2) (1.5,3) (2.5,0.3) (3,2.5) (4,3.5) (5,0.3) (5.5,0)};
  \node[font=\tiny, ufamgreen] at (3.8,1.0) {$R^{\delta}_I(J)$};
\end{tikzpicture}
\end{column}
\begin{column}{0.52\textwidth}
\textbf{Propriedades-chave:}
\begin{itemize}
  \item Propaga picos do marcador pelas regi\~oes conectadas da m\'ascara
  \item Para nas ``barreiras'' de intensidade
  \item \textbf{Preserva a forma} dos objetos remanescentes
  \item Superior \`a abertura padr\~ao: n\~ao encolhe bordas
\end{itemize}
\end{column}
\end{columns}

\end{frame}

%==============================================================
% SLIDE 8
%==============================================================
\begin{frame}{Abertura e Fechamento por Reconstru\c{c}\~ao}

\begin{columns}[T]
\begin{column}{0.48\textwidth}
\begin{block}{Abertura por Reconstru\c{c}\~ao}
\begin{enumerate}
  \item Erodir: $I_e = \varepsilon_B(I)$
  \item Reconstruir: $I_{\mathrm{obr}} = R^{\delta}_I(I_e)$
\end{enumerate}
\vspace{0.2cm}
\textbf{Efeito:} suprime picos claros (ru\'ido) \textbf{sem distorcer} as bordas reais dos objetos.
\end{block}
\end{column}
\begin{column}{0.48\textwidth}
\begin{block}{Fechamento por Reconstru\c{c}\~ao}
\begin{enumerate}
  \item Dilatar: $I_d = \delta_B(I)$
  \item Reconstruir no complemento
\end{enumerate}
\vspace{0.2cm}
\textbf{Efeito:} remove buracos escuros dentro dos objetos \textbf{preservando} sua geometria.
\end{block}
\end{column}
\end{columns}

\vspace{0.5cm}

\begin{exampleblock}{Vantagem sobre Abertura/Fechamento Padr\~ao}
A abertura padr\~ao $\gamma_B$ encolhe as bordas dos objetos. A abertura por reconstru\c{c}\~ao \textbf{restaura completamente} a forma de qualquer objeto que sobreviveu \`a eros\~ao --- apenas os picos menores que $B$ s\~ao eliminados.
\end{exampleblock}

\end{frame}

%==============================================================
\section{M\'etodo Proposto}
%==============================================================

%==============================================================
% SLIDE 9
%==============================================================
\begin{frame}{Vis\~ao Geral do Pipeline}

\begin{center}
\begin{tikzpicture}[
  node distance=0.3cm and 0.4cm,
  mybox/.style={draw, thick, rounded corners=3pt, text=white,
              font=\footnotesize\bfseries, minimum height=0.8cm,
              minimum width=2.0cm, align=center},
  arr/.style={->, thick, ufamblue},
  >=Stealth
]
  \node[mybox, fill=ufamblue] (input) {Imagem\\Original};
  \node[mybox, fill=ufamlightblue, right=of input] (grad) {Gradiente\\Morfol\'ogico};
  \node[mybox, fill=ufamgreen!80!black, right=of grad] (recon) {Abertura +\\Fechamento\\por Reconstr.};
  \node[mybox, fill=ufamgold!80!black, right=of recon, text=black] (fg) {Marcadores\\Foreground};
  \node[mybox, fill=ufamgold!80!black, below=0.5cm of fg, text=black] (bg) {Marcadores\\Background};
  \node[mybox, fill=alertred!80!black, right=of fg, yshift=-0.4cm] (impose) {Imposi\c{c}\~ao\\de M\'inimos};
  \node[mybox, fill=ufamblue, right=of impose] (ws) {Watershed\\Controlada};

  \draw[arr] (input) -- (grad);
  \draw[arr] (grad) -- (recon);
  \draw[arr] (recon) -- (fg);
  \draw[arr] (recon) |- (bg);
  \draw[arr] (fg) -- (impose);
  \draw[arr] (bg) -| (impose);
  \draw[arr] (impose) -- (ws);
\end{tikzpicture}
\end{center}

\vspace{0.3cm}

\begin{enumerate}
  \item \textbf{Gradiente morfol\'ogico} da imagem original
  \item \textbf{Limpeza} via abertura + fechamento por reconstru\c{c}\~ao (multi-escala)
  \item Extra\c{c}\~ao de \textbf{marcadores foreground} (m\'aximos regionais)
  \item Extra\c{c}\~ao de \textbf{marcadores background} (SKIZ via dist\^ancia)
  \item \textbf{Imposi\c{c}\~ao de m\'inimos} no gradiente
  \item \textbf{Watershed controlada por marcadores}
\end{enumerate}

\end{frame}

%==============================================================
% SLIDE 10
%==============================================================
\begin{frame}{Etapa 1 --- Gradiente Morfol\'ogico}

\begin{block}{C\'alculo}
\[
  g(I) = \delta_B(I) - \varepsilon_B(I)
\]
\end{block}

\begin{columns}[T]
\begin{column}{0.55\textwidth}
\textbf{Caracter\'isticas:}
\begin{itemize}
  \item Valores altos nas \textbf{transi\c{c}\~oes de borda}
  \item Valores $\approx 0$ em regi\~oes \textbf{homog\^eneas}
  \item Forma a superf\'icie topogr\'afica para a watershed
  \item Problema: cont\'em \textbf{muitos m\'inimos locais esp\'urios}
\end{itemize}

\vspace{0.3cm}
\textbf{Resultado:} cada m\'inimo local gerar\'a uma bacia de capta\c{c}\~ao $\Rightarrow$ sobre-segmenta\c{c}\~ao se usado diretamente.
\end{column}
\begin{column}{0.42\textwidth}
\centering
\begin{tikzpicture}[scale=0.55]
  % Simulated gradient with many minima
  \draw[thick, ufamblue] plot[smooth] coordinates {
    (0,1) (0.3,2.5) (0.6,0.8) (0.9,1.5) (1.2,0.3) (1.5,3.5)
    (1.8,0.5) (2.1,1.2) (2.4,0.4) (2.7,0.9) (3.0,3.8)
    (3.3,0.6) (3.6,1.0) (3.9,0.2) (4.2,2.8) (4.5,0.7)
    (4.8,1.3) (5.1,0.5) (5.4,1.8) (5.7,0.3)
  };
  \node[font=\tiny\bfseries] at (2.8,4.3) {Gradiente com m\'inimos esp\'urios};
  % Mark spurious minima
  \fill[alertred] (0.6,0.8) circle (2pt);
  \fill[alertred] (1.2,0.3) circle (2pt);
  \fill[alertred] (1.8,0.5) circle (2pt);
  \fill[alertred] (2.4,0.4) circle (2pt);
  \fill[alertred] (3.3,0.6) circle (2pt);
  \fill[alertred] (3.9,0.2) circle (2pt);
  \fill[alertred] (4.5,0.7) circle (2pt);
  \fill[alertred] (5.1,0.5) circle (2pt);
  \fill[alertred] (5.7,0.3) circle (2pt);
  \node[font=\tiny, alertred] at (3.0,-0.3) {$\bullet$ = m\'inimos locais};
\end{tikzpicture}
\end{column}
\end{columns}

\end{frame}

%==============================================================
% SLIDE 11
%==============================================================
\begin{frame}{Etapa 2 --- Eros\~ao Multi-Escala e Reconstru\c{c}\~ao}

\begin{block}{Inova\c{c}\~ao Central do Artigo}
Usar elementos estruturantes de \textbf{tamanhos diferentes} para cada escala de objeto:
\end{block}

\begin{columns}[T]
\begin{column}{0.48\textwidth}
\textbf{Abertura por Reconstru\c{c}\~ao:}
\begin{enumerate}
  \item $I_e = \varepsilon_{B_k}(I)$ com $B_k$ adaptado \`a escala $k$
  \item $I_{\mathrm{obr}} = R^{\delta}_I(I_e)$
\end{enumerate}
\vspace{0.2cm}
$B$ pequeno $\to$ preserva detalhes finos\\
$B$ grande $\to$ elimina ru\'ido em regi\~oes grandes

\vspace{0.3cm}
\textbf{Fechamento por Reconstru\c{c}\~ao:}
\begin{enumerate}
  \item $I_d = \delta_{B_k}(I_{\mathrm{obr}})$
  \item Reconstru\c{c}\~ao no complemento
\end{enumerate}
\end{column}
\begin{column}{0.48\textwidth}
\centering
\begin{tikzpicture}[scale=0.7]
  % Multi-scale SE illustration
  \fill[blue!20] (0,0) rectangle (0.4,0.4);
  \draw[thick] (0,0) rectangle (0.4,0.4);
  \node[font=\tiny, below] at (0.2,-0.1) {$B_1$ (3$\times$3)};

  \fill[blue!30] (1.5,-0.2) rectangle (2.3,0.6);
  \draw[thick] (1.5,-0.2) rectangle (2.3,0.6);
  \node[font=\tiny, below] at (1.9,-0.3) {$B_2$ (5$\times$5)};

  \fill[blue!40] (3.2,-0.5) rectangle (4.4,0.7);
  \draw[thick] (3.2,-0.5) rectangle (4.4,0.7);
  \node[font=\tiny, below] at (3.8,-0.6) {$B_3$ (9$\times$9)};

  % Objects they target
  \draw[->, thick, ufamgreen] (0.2,0.6) -- (0.2,1.3);
  \node[font=\tiny, ufamgreen, above] at (0.2,1.3) {Obj.\ peq.};

  \draw[->, thick, ufamgold!80!black] (1.9,0.8) -- (1.9,1.3);
  \node[font=\tiny, ufamgold!80!black, above] at (1.9,1.3) {Obj.\ m\'ed.};

  \draw[->, thick, ufamblue] (3.8,0.9) -- (3.8,1.3);
  \node[font=\tiny, ufamblue, above] at (3.8,1.3) {Obj.\ grande};
\end{tikzpicture}

\vspace{0.4cm}
{\footnotesize As bacias de capta\c{c}\~ao s\~ao semeadas separadamente em m\'ultiplas escalas.}
\end{column}
\end{columns}

\end{frame}

%==============================================================
% SLIDE 12
%==============================================================
\begin{frame}{Etapa 3 --- Extra\c{c}\~ao de Marcadores Foreground}

\begin{block}{Procedimento}
\begin{enumerate}
  \item Extrair os \textbf{m\'aximos regionais} da imagem reconstru\'ida:
  \[
    M_{fg} = \texttt{imregionalmax}(I_{\mathrm{obrcbr}})
  \]
  \item Remover blobs muito pequenos (\'area $<$ limiar) com \texttt{bwareaopen}
  \item Resultado: cada componente conexa \'e um \textbf{marcador foreground} (semente de bacia)
\end{enumerate}
\end{block}

\vspace{0.3cm}

\begin{exampleblock}{M\'aximos Regionais}
Uma componente conexa de pixels \'e um \textbf{m\'aximo regional} se todos os seus vizinhos t\^em valor \textbf{estritamente menor}. Ap\'os a reconstru\c{c}\~ao, os objetos t\^em interiores planos, gerando m\'aximos regionais bem definidos exatamente dentro de cada objeto.
\end{exampleblock}

\end{frame}

%==============================================================
% SLIDE 13
%==============================================================
\begin{frame}{Etapa 4 --- Extra\c{c}\~ao de Marcadores Background (SKIZ)}

\begin{block}{Skeleton by Influence Zones (SKIZ)}
\begin{enumerate}
  \item Binarizar a imagem limpa $\to$ m\'ascara foreground
  \item Calcular a \textbf{transformada de dist\^ancia} $D$ do background
  \item Aplicar watershed a $D$: as linhas de crista ($DL = 0$) s\~ao os marcadores background
\end{enumerate}
\end{block}

\begin{center}
\begin{tikzpicture}[scale=0.7, >=Stealth]
  % Objects
  \fill[ufamblue!30] (0.5,0.5) circle (0.6);
  \fill[ufamblue!30] (3.0,1.0) circle (0.9);
  \fill[ufamblue!30] (1.5,2.5) circle (0.5);

  % SKIZ lines
  \draw[thick, alertred, dashed] (1.5,0) -- (1.8,1.5) -- (0.5,2.0) -- (0,2.5);
  \draw[thick, alertred, dashed] (1.8,1.5) -- (2.5,2.8) -- (3.5,3.0);
  \draw[thick, alertred, dashed] (1.8,1.5) -- (4.0,0);

  % Labels
  \node[font=\tiny\bfseries, ufamblue] at (0.5,0.5) {$O_1$};
  \node[font=\tiny\bfseries, ufamblue] at (3.0,1.0) {$O_2$};
  \node[font=\tiny\bfseries, ufamblue] at (1.5,2.5) {$O_3$};
  \node[font=\tiny, alertred] at (3.5,2.2) {SKIZ (bg markers)};

  \draw[thick] (0,0) rectangle (4.0,3.0);
\end{tikzpicture}
\end{center}

{\footnotesize Os marcadores background ficam no ``ponto m\'edio'' entre objetos, evitando que as linhas de watershed invadam os objetos.}

\end{frame}

%==============================================================
% SLIDE 14
%==============================================================
\begin{frame}{Etapa 5 --- Imposi\c{c}\~ao de M\'inimos}

\begin{block}{Modifica\c{c}\~ao de Homotopia}
\[
  G_{\mathrm{mod}} = \texttt{imimposemin}\!\left(G,\; M_{fg} \cup M_{bg}\right)
\]
\end{block}

\begin{columns}[T]
\begin{column}{0.55\textwidth}
\textbf{O que faz:}
\begin{itemize}
  \item Reescreve o gradiente para que m\'inimos regionais existam \textbf{somente} nas posi\c{c}\~oes dos marcadores
  \item Todos os outros m\'inimos locais s\~ao ``preenchidos'' at\'e o n\'ivel vizinho
  \item Preserva a rela\c{c}\~ao de ordem relativa
  \item Implementado via reconstru\c{c}\~ao geod\'esica no complemento
\end{itemize}
\end{column}
\begin{column}{0.42\textwidth}
\centering
\begin{tikzpicture}[scale=0.55]
  % Before
  \node[font=\tiny\bfseries, ufamblue] at (2.5,4.5) {Antes};
  \draw[thick, gray] plot[smooth] coordinates {
    (0,2) (0.5,1) (1,2.5) (1.5,0.5) (2,3) (2.5,0.3) (3,2.8) (3.5,1.2) (4,2) (4.5,0.8) (5,2)
  };
  \fill[alertred] (0.5,1) circle (2pt);
  \fill[alertred] (1.5,0.5) circle (2pt);
  \fill[alertred] (2.5,0.3) circle (2pt);
  \fill[alertred] (3.5,1.2) circle (2pt);
  \fill[alertred] (4.5,0.8) circle (2pt);
\end{tikzpicture}

\vspace{0.2cm}

\begin{tikzpicture}[scale=0.55]
  % After
  \node[font=\tiny\bfseries, ufamgreen] at (2.5,4.5) {Depois};
  \draw[thick, ufamgreen!80!black] plot[smooth] coordinates {
    (0,2) (0.8,1.8) (1,2.5) (1.5,1.8) (2,3) (2.5,0.3) (3,2.8) (3.2,2.5) (4,2) (4.3,1.8) (5,2)
  };
  \fill[ufamgreen] (2.5,0.3) circle (3pt);
  \node[font=\tiny, ufamgreen] at (2.5,-0.3) {Marcador};
\end{tikzpicture}
\end{column}
\end{columns}

\end{frame}

%==============================================================
% SLIDE 15
%==============================================================
\begin{frame}{Etapa 6 --- Watershed Controlada por Marcadores}

\begin{block}{Resultado Final}
A watershed \'e aplicada ao gradiente modificado $G_{\mathrm{mod}}$:
\begin{itemize}
  \item Todos os m\'inimos esp\'urios foram eliminados
  \item A inunda\c{c}\~ao come\c{c}a \textbf{somente} nos marcadores definidos
  \item Cada bacia de capta\c{c}\~ao corresponde a \textbf{exatamente um objeto}
  \item Linhas de watershed coincidem com \textbf{bordas reais}
\end{itemize}
\end{block}

\vspace{0.3cm}

\begin{center}
\begin{tikzpicture}[scale=0.8, >=Stealth]
  % Traditional watershed (over-segmented)
  \node[font=\small\bfseries, alertred] at (2,3.5) {Watershed Tradicional};
  \draw[thick] (0,0) rectangle (4,3);
  % Many small regions
  \draw[thin, red!50] (0.5,0) -- (0.5,3);
  \draw[thin, red!50] (1.0,0) -- (1.0,3);
  \draw[thin, red!50] (1.5,0) -- (1.5,3);
  \draw[thin, red!50] (2.0,0) -- (2.0,3);
  \draw[thin, red!50] (2.5,0) -- (2.5,3);
  \draw[thin, red!50] (3.0,0) -- (3.0,3);
  \draw[thin, red!50] (3.5,0) -- (3.5,3);
  \draw[thin, red!50] (0,0.5) -- (4,0.5);
  \draw[thin, red!50] (0,1.0) -- (4,1.0);
  \draw[thin, red!50] (0,1.5) -- (4,1.5);
  \draw[thin, red!50] (0,2.0) -- (4,2.0);
  \draw[thin, red!50] (0,2.5) -- (4,2.5);
  \node[font=\footnotesize, alertred] at (2,1.5) {Sobre-segmentado};

  % Proposed method
  \begin{scope}[xshift=6cm]
  \node[font=\small\bfseries, ufamgreen] at (2,3.5) {M\'etodo Proposto};
  \draw[thick] (0,0) rectangle (4,3);
  % Clean regions
  \draw[ultra thick, ufamgreen] (0,0) -- (0,3) -- (4,3) -- (4,0) -- cycle;
  \draw[ultra thick, ufamgreen] plot[smooth cycle] coordinates {(0.5,0.5) (1.5,0.3) (2.0,1.0) (1.5,1.5) (0.5,1.2)};
  \draw[ultra thick, ufamgreen] plot[smooth cycle] coordinates {(2.5,1.5) (3.5,1.3) (3.7,2.5) (2.8,2.7)};
  \draw[ultra thick, ufamgreen] (0.3,2.0) circle (0.4);
  \node[font=\footnotesize, ufamgreen] at (2,-0.4) {Segmenta\c{c}\~ao precisa};
  \end{scope}

  \draw[->, ultra thick, ufamblue] (4.3,1.5) -- (5.7,1.5);
\end{tikzpicture}
\end{center}

\end{frame}

%==============================================================
\section{Formul\c{c}\~ao Matem\'atica}
%==============================================================

%==============================================================
% SLIDE 16
%==============================================================
\begin{frame}{Resumo Matem\'atico do Algoritmo}

\begin{enumerate}
  \item \textbf{Gradiente morfol\'ogico:}
  \[
    G = \delta_B(I) - \varepsilon_B(I)
  \]

  \item \textbf{Abertura por reconstru\c{c}\~ao:}
  \[
    I_{\mathrm{obr}} = R^{\delta}_I\!\left(\varepsilon_{B_k}(I)\right)
  \]

  \item \textbf{Fechamento por reconstru\c{c}\~ao:}
  \[
    I_{\mathrm{obrcbr}} = \left[R^{\varepsilon}_{I_{\mathrm{obr}}}\!\left(\delta_{B_k}(I_{\mathrm{obr}})\right)\right]^c
  \]

  \item \textbf{Marcadores foreground:}
  $M_{fg} = \texttt{imregionalmax}(I_{\mathrm{obrcbr}})$

  \item \textbf{Marcadores background (SKIZ):}
  $M_{bg} = \texttt{watershed}\!\left(\texttt{bwdist}(\overline{M_{fg}})\right) == 0$

  \item \textbf{Imposi\c{c}\~ao:}
  $G_{\mathrm{mod}} = \texttt{imimposemin}(G,\; M_{fg} \cup M_{bg})$

  \item \textbf{Segmenta\c{c}\~ao final:}
  $L = \texttt{watershed}(G_{\mathrm{mod}})$
\end{enumerate}

\end{frame}

%==============================================================
\section{Resultados Experimentais}
%==============================================================

%==============================================================
% SLIDE 17
%==============================================================
\begin{frame}{Resultados Experimentais}

\begin{columns}[T]
\begin{column}{0.48\textwidth}
\begin{block}{Configura\c{c}\~ao}
\begin{itemize}
  \item Imagens complexas com objetos de \textbf{m\'ultiplos tamanhos}
  \item Compara\c{c}\~ao com:
  \begin{itemize}
    \item Watershed tradicional (direto no gradiente)
    \item M\'etodos de limiar \'unico
    \item Crescimento de regi\~oes
    \item Limiariza\c{c}\~ao OTSU
  \end{itemize}
\end{itemize}
\end{block}
\end{column}
\begin{column}{0.48\textwidth}
\begin{exampleblock}{Resultados Quantitativos}
\renewcommand{\arraystretch}{1.3}
\begin{tabular}{lc}
\toprule
\textbf{M\'etrica} & \textbf{Valor} \\
\midrule
Precis\~ao & $> 0{,}98$ \\
Recall & $> 0{,}98$ \\
Melhoria de efici\^encia & $\sim 10\%$ \\
\bottomrule
\end{tabular}
\end{exampleblock}
\end{column}
\end{columns}

\vspace{0.4cm}

\begin{alertblock}{Principais Achados}
\begin{itemize}
  \item M\'etodos de limiar \'unico ou fragmentam objetos grandes ou perdem objetos pequenos
  \item O m\'etodo proposto delineia \textbf{ambos simultaneamente} em uma \'unica passada
  \item Contornos fechados, com \textbf{largura de 1 pixel}, m\'inima sub/sobre-segmenta\c{c}\~ao
  \item \textbf{Robusto} a ru\'ido salt-and-pepper
\end{itemize}
\end{alertblock}

\end{frame}

%==============================================================
% SLIDE 18
%==============================================================
\begin{frame}{Compara\c{c}\~ao Visual de Resultados}

\begin{center}
\begin{tikzpicture}[scale=0.75, >=Stealth]
  % Column 1 - Original
  \node[font=\small\bfseries, ufamblue] at (1.5,4.5) {Original};
  \draw[thick] (0,1) rectangle (3,4);
  % Simple objects
  \fill[gray!60] (0.5,1.5) circle (0.3);
  \fill[gray!40] (1.5,2.5) ellipse (0.8 and 0.6);
  \fill[gray!50] (2.5,3.5) circle (0.2);

  % Column 2 - Traditional Watershed
  \begin{scope}[xshift=4cm]
  \node[font=\small\bfseries, alertred] at (1.5,4.5) {Watershed Direta};
  \draw[thick] (0,1) rectangle (3,4);
  \fill[gray!60] (0.5,1.5) circle (0.3);
  \fill[gray!40] (1.5,2.5) ellipse (0.8 and 0.6);
  \fill[gray!50] (2.5,3.5) circle (0.2);
  % Over-segmentation lines
  \draw[red, thick] (0,2.0) -- (3,2.0);
  \draw[red, thick] (0,3.0) -- (3,3.0);
  \draw[red, thick] (1.0,1) -- (1.0,4);
  \draw[red, thick] (2.0,1) -- (2.0,4);
  \draw[red, thick] (0.5,1) -- (2.5,4);
  \node[font=\tiny, alertred] at (1.5,0.7) {$\sim$200 regi\~oes};
  \end{scope}

  % Column 3 - Proposed method
  \begin{scope}[xshift=8cm]
  \node[font=\small\bfseries, ufamgreen] at (1.5,4.5) {M\'etodo Proposto};
  \draw[thick] (0,1) rectangle (3,4);
  \draw[ultra thick, ufamgreen] (0.5,1.5) circle (0.3);
  \draw[ultra thick, ufamgreen] (1.5,2.5) ellipse (0.8 and 0.6);
  \draw[ultra thick, ufamgreen] (2.5,3.5) circle (0.2);
  \node[font=\tiny, ufamgreen] at (1.5,0.7) {3 regi\~oes (correto)};
  \end{scope}
\end{tikzpicture}
\end{center}

\vspace{0.3cm}

\begin{columns}[T]
\begin{column}{0.48\textwidth}
{\footnotesize\textbf{Watershed direta:} Cada m\'inimo esp\'urio do gradiente gera uma regi\~ao separada, resultando em fragmenta\c{c}\~ao excessiva.}
\end{column}
\begin{column}{0.48\textwidth}
{\footnotesize\textbf{M\'etodo proposto:} Marcadores multi-escala + reconstru\c{c}\~ao produzem uma bacia por objeto, com contornos precisos.}
\end{column}
\end{columns}

\end{frame}

%==============================================================
\section{Conclus\~oes}
%==============================================================

%==============================================================
% SLIDE 19
%==============================================================
\begin{frame}{Contribui\c{c}\~oes do Artigo}

\begin{enumerate}
  \item \textbf{Eros\~ao multi-escala para gera\c{c}\~ao de marcadores}
  \begin{itemize}
    \item Elementos estruturantes adaptados ao tamanho de cada objeto
    \item Semeia bacias de capta\c{c}\~ao em todas as escalas simultaneamente
  \end{itemize}

  \vspace{0.3cm}

  \item \textbf{Reconstru\c{c}\~ao morfol\'ogica para limpeza do gradiente}
  \begin{itemize}
    \item Superior \`a abertura/fechamento padr\~ao: preserva bordas
    \item Remove ru\'ido sem distorcer a geometria dos objetos
  \end{itemize}

  \vspace{0.3cm}

  \item \textbf{Extra\c{c}\~ao sistem\'atica de marcadores FG + BG}
  \begin{itemize}
    \item M\'aximos regionais (foreground) + SKIZ (background)
    \item Garante que cada linha de watershed cai sobre uma borda real
  \end{itemize}

  \vspace{0.3cm}

  \item \textbf{Pipeline integrado}
  \begin{itemize}
    \item Segmenta\c{c}\~ao com contornos fechados e precisos em imagens multi-escala
    \item Aplic\'avel a imagens biom\'edicas, sensoriamento remoto, inspe\c{c}\~ao industrial
  \end{itemize}
\end{enumerate}

\end{frame}

%==============================================================
% SLIDE 20
%==============================================================
\begin{frame}{Conclus\~ao e Mensagem Principal}

\begin{block}{Mensagem Principal}
A \textbf{qualidade dos marcadores} --- e n\~ao a transformada watershed em si --- \'e o fator determinante da qualidade da segmenta\c{c}\~ao. O m\'etodo proposto dedica o n\'ucleo do algoritmo \`a computa\c{c}\~ao rigorosa desses marcadores via reconstru\c{c}\~ao morfol\'ogica multi-escala.
\end{block}

\vspace{0.5cm}

\begin{exampleblock}{Aplica\c{c}\~oes}
\begin{columns}[T]
\begin{column}{0.3\textwidth}
\centering
\textbf{Biom\'edica}\\
C\'elulas, tecidos, \'org\~aos
\end{column}
\begin{column}{0.3\textwidth}
\centering
\textbf{Sensoriamento Remoto}\\
Edif\'icios, vegeta\c{c}\~ao
\end{column}
\begin{column}{0.3\textwidth}
\centering
\textbf{Inspe\c{c}\~ao Industrial}\\
Defeitos, pe\c{c}as
\end{column}
\end{columns}
\end{exampleblock}

\vspace{0.5cm}

\begin{center}
{\Large\bfseries\color{ufamblue} Obrigado!}\\[0.3cm]
{\small Wu, Y., Peng, X., Ruan, K. \textit{et al.} \textit{Multimedia Tools and Applications}, vol.\ 76, pp.\ 19781--19793, 2017.}
\end{center}

\end{frame}

%==============================================================
% SLIDE 21 — REFERÊNCIAS
%==============================================================
\begin{frame}{Refer\^encias}

\begin{thebibliography}{9}
\footnotesize

\bibitem{wu2017}
Wu, Y., Peng, X., Ruan, K. \textit{et al.}
\textit{Improved image segmentation method based on morphological reconstruction}.
Multimedia Tools and Applications, v.\ 76, p.\ 19781--19793, 2017.

\bibitem{gonzalez2018}
Gonzalez, R.\ C.; Woods, R.\ E.
\textit{Digital Image Processing}. 4.\ ed.\ Pearson, 2018.

\bibitem{soille2003}
Soille, P.
\textit{Morphological Image Analysis: Principles and Applications}. 2.\ ed.\ Springer, 2003.

\bibitem{vincent1993}
Vincent, L.; Soille, P.
Watersheds in digital spaces: an efficient algorithm based on immersion simulations.
\textit{IEEE Trans.\ PAMI}, v.\ 13, n.\ 6, p.\ 583--598, 1991.

\bibitem{matlab}
MathWorks.
\textit{Marker-Controlled Watershed Segmentation}.
MATLAB Documentation, 2024.

\end{thebibliography}

\end{frame}

\end{document}
